\chapter{Literature Review}
\label{chap:lit}

\section{Survival Data and the Prediction Task}
\label{survival-analysis-and-right-censored-data}

Survival methods connect an event definition, an observation process and a model for time to that event. In and Lee explain how event times and risk sets define survival estimates, while Perera and Dwivedi relate these quantities to study design \cite{ravi14,ravi15}. Their central implication for prediction is that the endpoint and time origin define the task before a learning algorithm is selected. Mortality and recurrence models can use similar statistical machinery while answering different clinical questions.

Censoring also links estimation to evaluation. Survival likelihoods and inverse-probability-weighted scores use partial follow-up under assumptions about the censoring process \cite{ravi15,ravi26}. At a fixed horizon, early-censored patients have unknown outcome status; changing the algorithm does not remove that observation problem. A suitable comparison therefore aligns endpoints, prediction horizons and censoring adjustment as well as predictor inputs.

Observed events, follow-up and predictor dependence jointly determine the information available to estimate effects. Comparing models across clinical cohorts retains real differences in these quantities, whereas matched changes within a cohort examine information more directly. These two designs provide complementary evidence about prior behaviour. The event-time notation and the implemented censoring rules are specified in Chapter~\ref{chap:methods}.

\section{Regression and Survival Ensembles}
\label{classical-survival-models}

\subsection{Descriptive Survival Estimates and Cox Regression}

The Kaplan--Meier estimator describes the survival function through successive risk sets, while the Nelson--Aalen estimator describes cumulative hazard. These methods show how survival changes over time without requiring a multivariable regression structure. Group comparisons such as the log-rank test evaluate differences in observed survival experience. They provide a descriptive context for the clinical data alongside patient-level predictive models \cite{ravi14,ravi16,ravi18}.

Cox regression specifies
\[
 h(t\mid x)=h_0(t)\exp(x^\top\beta),
\]
where $h_0(t)$ is the baseline hazard and $\beta$ contains the predictor coefficients. The partial likelihood estimates relative effects using risk sets without specifying a parametric distribution for the baseline hazard \cite{cox1972}. Absolute survival prediction additionally uses an estimate of the cumulative baseline hazard. A coefficient has a direct conditional interpretation: changing one encoded predictor unit multiplies the hazard by $\exp(\beta_j)$, with the other predictors held fixed.

Proportional hazards and the form of the predictor are separate modelling choices. Proportional hazards makes a covariate contrast constant over time. A linear predictor assigns a constant log-hazard change per predictor unit; transformations and interactions can represent other relationships. Jachno et al. discuss the practical importance of recognising non-constant hazards and non-proportional effects in clinical trials \cite{ravi17}. Their review connects model assumptions to the scientific question rather than treating a familiar model as sufficient justification for every event pattern.

CoxPH is a useful comparator because its prediction structure is clear and its coefficients support clinical interpretation. It also provides a close reference for a Bayesian proportional-hazards model: both relate predictors to hazard multiplicatively, while differing in likelihood treatment, regularisation and uncertainty representation. These differences help explain why comparisons among the Bayesian prior variants are more direct measures of prior choice than comparisons between Bayesian and conventional Cox estimation.

\subsection{Random Survival Forests and Other Learning Objectives}
\label{machine-learning-survival-models}

Random survival forest extends the random-forest approach to right-censored outcomes. Survival trees are fitted to bootstrap samples, random subsets of predictors are considered at splits, and tree-level survival or cumulative-hazard estimates are aggregated \cite{ishwaran2008}. The resulting ensemble represents nonlinear relationships and interactions through partitions of the predictor space. Its event-time predictions arise from the survival experience within those partitions rather than a single linear proportional-hazards risk function.

The internal design of a forest matters. Schmid et al. examine concordance-based splitting as an alternative to log-rank splitting \cite{ravi19}. This comparison shows how the objective used to form a tree can change the prediction task emphasised during training. Tree depth, terminal-node size and the number of candidate predictors also determine the balance between local flexibility and stable estimation. RSF should therefore be understood as a specified learning procedure, including its tuning rules, rather than a label that uniquely determines predictive behaviour.

Other survival-learning methods illustrate the same distinction between model structure and objective. Fouodo et al. compare regression, ranking and hybrid support-vector formulations with established survival approaches \cite{ravi20}. Ranking methods focus on ordering comparable patients; regression formulations target survival-time relationships under censoring constraints. Boosting procedures similarly depend on their chosen loss and base learner. In the Fouodo comparison, the boosting reference used a proportional-hazards base learner. That result concerns a particular implementation, rather than all possible tree-boosting methods.

This literature supports comparing methods through a shared evaluation design. A flexible model can learn a useful interaction, while a more constrained model can estimate a stable relationship from fewer events. Simulation comparisons of random forests and Cox regression, and reviews of forests for high-dimensional data, examine these interactions between model specification and data structure \cite{ravi37,ravi52}. The relevant empirical question is how a specified procedure behaves on held-out observations under common outcome definitions and measures.

Interpretation also changes across these models. An RSF permutation score summarises how a selected prediction measure changes after perturbing a feature. A regression coefficient describes a conditional association on an encoded log-hazard scale. The forest importance distribution studied by Ishwaran and Lu provides a statistical treatment of importance uncertainty \cite{ravi70}. It supports interpreting importance together with its definition and variability. In this thesis, RSF supplies a nonlinear ensemble comparator for prediction, while posterior quantities provide the main distributional interpretation of the Bayesian models.

\section{Neural Survival Models}
\label{deep-learning-survival-models}

Neural survival models learn a predictor representation while using an objective adapted to censoring and event timing. Their outputs differ: some estimate a risk function, some predict a discrete event-time distribution, and others parameterise a continuous-time distribution. Wang et al. review this wider landscape, while Kvamme and Borgan compare continuous- and discrete-time neural formulations \cite{ravi3,ravi21}. The representation of time is therefore part of the model, alongside the architecture of the network.

DeepSurv replaces the linear Cox risk predictor with a feed-forward neural network and trains it using the Cox partial likelihood \cite{ravi4}. A nonlinear log-risk function can represent associations and interactions that would otherwise need explicit regression terms. The proportional-hazards structure remains: a covariate contrast changes relative risk through the learned function, while the shared baseline describes time. This makes DeepSurv a relevant comparator when the Bayesian model also uses proportional hazards but retains a linear predictor.

The original DeepSurv experiments include clinical datasets related to those used here. Their relevance lies in the modelling approach and clinical context. Numerical comparisons depend on the exact cohort version, feature representation, training procedure and evaluation design. A common set of folds and prediction horizons places DeepSurv, CoxPH and the Bayesian models on the same empirical task in the present study.

DeepHit offers a different construction. It learns a discrete event-time probability distribution using likelihood and ranking components, with a framework that accommodates competing risks \cite{lee2018deephit}. Kvamme and Borgan show how continuous- and discrete-time formulations create different modelling and estimation choices \cite{ravi21}. A discrete grid controls how event times are represented, whereas a proportional-hazards risk function separates predictor effects from the baseline hazard. These methods demonstrate that neural flexibility can be introduced into the predictor, the time distribution or both.

Kantidakis, Hazewinkel and Fiocco review neural survival prediction in medicine and critically appraise its use with prognostic factors \cite{ravi5}. The reviewed studies vary in endpoint, cohort size, architecture and validation. Their collective value is the range of modelling strategies and evaluation practices they reveal. Survival benchmarks add comparisons under shared procedures, connecting that methodological diversity to reproducible empirical assessment \cite{ravi33,ravi79}.

For the present comparison, DeepSurv represents nonlinear neural risk prediction within the PH setting. CoxPH represents conventional regression, and RSF represents a survival ensemble. This set provides several forms of predictive flexibility around the central comparison of Bayesian coefficient priors. The distinction between model-family comparisons and controlled prior comparisons remains clear because only the latter share the same likelihood and baseline-hazard model.

\section{Bayesian Survival Models and Baseline Hazards}
\label{bayesian-survival-modelling}

Bayesian inference combines a likelihood with a prior distribution to obtain a posterior distribution. For observed data $D$ and parameters $\theta$,
\[
 p(\theta\mid D)\propto p(D\mid\theta)p(\theta).
\]
In survival analysis, the likelihood includes observed events and event-free exposure to censoring. Posterior draws then describe uncertainty in coefficients and in model-derived predictions. Alvares et al. present Bayesian formulations for proportional hazards, accelerated failure time, frailty, competing risks and other event structures \cite{ravi22}. Bayesian survival analysis is consequently a framework for several statistical models, with the likelihood chosen to match the outcome and scientific question.

\subsection{Full Likelihood and Piecewise Baseline Hazards}

A full-likelihood proportional-hazards model assigns a structure to both predictor effects and the baseline hazard. A piecewise-exponential specification divides follow-up into intervals and assigns a constant baseline rate within each interval. The cumulative hazard is the sum of interval rates multiplied by the time exposed within them. Friedman, and Laird and Olivier, establish the relationship between this representation and event/exposure regression \cite{friedman1982,laird1981}. It provides a tractable route from censored observations to coefficient and baseline-hazard estimation.

Interval-specific rates allow the baseline to change over time while retaining a common covariate risk function. L\'{a}zaro, Armero and Alvares study regularisation of flexible baseline hazards in Bayesian Cox-type survival models \cite{ravi23}. Their work shows how baseline flexibility and prior structure can be considered together. More intervals create more local rates to estimate, while regularisation influences how much those estimates can vary. Holding the interval construction and baseline prior constant creates a controlled setting for evaluating coefficient priors.

Qing, Thall and Yuan apply a Bayesian piecewise-exponential model to monitoring a time-to-event endpoint in a phase II trial \cite{ravi38}. The application demonstrates the usefulness of a time-varying baseline representation for a specific decision problem. Freeman et al. compare piecewise-exponential and other flexible formulations in Bayesian meta-analysis with non-proportional hazards \cite{ravi96}. These studies concern trial monitoring and evidence synthesis, respectively. Their shared methodological relevance is the explicit representation of time and the importance of aligning that representation with the estimand.

In the present study, six training-derived baseline intervals define the primary Bayesian model, with other interval counts examined as sensitivities. The baseline can change between intervals, while the coefficient effects retain proportional hazards. This separation gives the coefficient-prior comparison a clear target: different regularisation of the same predictor effects within one survival likelihood.

\subsection{Posterior Computation and Diagnostic Assessment}

Posterior inference may use sampling or deterministic approximation. Alvares et al. describe survival modelling with Integrated Nested Laplace Approximation, which exploits latent Gaussian structure for approximate Bayesian inference \cite{ravi94}. NUTS instead extends Hamiltonian Monte Carlo by adapting trajectory length during sampling \cite{hoffman2014}. The choice of inference method interacts with the likelihood, parameterisation and prior hierarchy.

Hierarchical priors can create posterior regions with different scales. Computational assessment is therefore part of evaluating a prior specification. Rank-normalised $\widehat R$, effective sample size, divergent transitions, energy behaviour and tree-depth diagnostics examine different features of finite-chain sampling \cite{ravi39,standiagnostics}. These quantities complement posterior estimates by describing how the draws were obtained. In this study, the same diagnostic acceptance criteria are applied across Bayesian arms, linking prior comparison to both predictive output and computational reliability.

\section{Coefficient Priors, Shrinkage and Information}
\label{bayesian-priors-and-shrinkage}

A coefficient prior defines a distribution of plausible effects on the model's predictor scale. Standardisation makes a common coefficient scale interpretable across numeric variables, while categorical encoding determines the contrasts assigned separate coefficients. The prior operates on that encoded representation. Its interaction with the likelihood depends on event information and on the degree to which predictors provide distinct evidence about their effects \cite{ibrahim2001,ravi22}.

\subsection{Normal and Global--Local Shrinkage}

An independent zero-centred Normal prior provides a common scale for regression coefficients. It regularises large values while retaining a continuous range of effects. Its role as a reference is useful because the assumed coefficient scale is explicit. Comparing other priors against this reference measures what a more adaptive scale structure changes under the same likelihood.

The horseshoe prior combines a global scale, which controls overall shrinkage, with local scales that allow individual coefficients to respond differently \cite{carvalho2010}. A coefficient with little data support can be concentrated near zero, while a coefficient with stronger support can retain a larger magnitude. Bhadra et al. develop the related Horseshoe+ estimator for sparse signals, illustrating how hierarchical scale distributions create different shrinkage behaviour \cite{ravi24}.

Piironen and Vehtari introduce finite-slab regularisation to the horseshoe and relate global-scale choice to prior information about sparsity \cite{piironen2017}. The slab controls the largest effective coefficient scales. This provides a useful distinction between concentrating small effects and controlling large effects: both contribute to the resulting prior distribution. The specification used in this thesis follows the regularised global--local construction with half-$t_2$ scale distributions. The exact distributions and numerical scales are defined in Chapter~\ref{chap:methods}, making the fitted model explicit alongside the broader horseshoe literature.

The computational consequences are also relevant. A hierarchical prior estimates more scale parameters than an independent Normal prior. Its posterior can contain strong scale dependence, especially where effects and their local scales are weakly separated by the data. The shared diagnostic assessment therefore contributes information about the behaviour of the fitted prior under each cohort's event information.

\subsection{Spike-and-Slab Priors and Continuous Mixtures}

Spike-and-slab approaches distinguish a narrow coefficient distribution from a broader one. In discrete variable-selection formulations, an inclusion indicator determines the component assigned to a coefficient \cite{george1993}. Tang et al. apply a spike-and-slab lasso Cox model to survival prediction and associated-gene detection, while Samorodnitsky and colleagues develop hierarchical spike-and-slab modelling for pan-cancer survival with pan-omic data \cite{ravi25,ravi43}. These studies connect prior structure to shrinkage and feature selection in survival settings.

Nikooienejad, Wang and Johnson use a point mass at zero and a non-local inverse moment prior with a Cox partial likelihood \cite{ravi6}. Their construction emphasises model selection in high-dimensional survival data. Together, these approaches show that the distribution assigned to small effects, the treatment of larger effects and the model-space prior jointly determine selection behaviour. The word ``spike-and-slab'' encompasses several statistical constructions, so the implemented distribution is essential to interpretation.

This thesis uses a continuous variance interpolation. A Beta-distributed weight moves a coefficient scale between narrow and broad values, and the coefficient itself remains continuously distributed. The resulting weight describes that scale interpolation. Posterior summaries of the weight accompany coefficient estimates as prior-specific information, while posterior expected contributions provide the common importance measure across all three priors. This choice gives a coherent comparison despite differences in their internal parameters.

\subsection{Training Information and Prior Sensitivity}

The likelihood and prior jointly determine the posterior. When event information is strong and predictors are well separated, coefficient estimates can be similar across reasonable priors. Fewer observed events or overlapping predictor information can make prior structure more influential. Studies of prior impact provide ways to connect prior assumptions to effective information, while Bayesian survival applications illustrate how the chosen likelihood and coefficient structure shape the resulting inference \cite{ravi57,ravi68,ravi98}.

Training size and censoring affect information differently. A smaller training sample reduces the observations available for estimation. Shorter follow-up changes which events are observed and how much time patients contribute. Controlled changes to both provide a more informative assessment than comparing cohorts only by sample size. The present study uses shared training subsets and training-derived administrative truncation so that prior effects can be examined under matched information conditions.

Sensitivity also concerns design choices that alter the representation of the same clinical data. Alternative baseline intervals, predictor specifications, missing-data procedures and coefficient scales can change both estimation and prediction. Evaluating these conditions within the same validation structure connects a prior's headline performance to the settings in which it is used. The resulting evidence concerns the defined prior implementations and clinical tasks, with computational diagnostics retained as part of each comparison.

\section{Evaluating Survival Predictions}
\label{survival-model-evaluation}

Survival prediction has several dimensions of performance. Discrimination concerns patient ordering, probability error concerns the numerical survival predictions, and calibration concerns agreement between predicted and observed event probabilities. Park et al. review these distinctions for conventional and deep-learning survival methods, while McLernon et al. provide practical guidance for evaluating predictions with survival outcomes \cite{ravi26,ravi88}. A common set of complementary measures gives a more complete empirical description than a single model ranking.

\subsection{Discrimination}

Harrell's concordance summarises agreement between predicted risk and observed event ordering among comparable patient pairs \cite{harrell1982}. Uno's estimator uses censoring weights to address the censoring dependence of conventional concordance summaries \cite{uno2011}. Longato et al. explain the practical meaning of concordance and its role in model evaluation \cite{ravi27}. Both estimators concern relative ordering rather than the absolute value of a survival probability.

Time-dependent AUC examines discrimination at a chosen horizon \cite{heagerty2000,ravi26}. Under a cumulative/dynamic definition, cases have experienced the event by time $t$ and controls remain event-free beyond $t$. Its time dependence is useful when a model separates risks more clearly during one part of follow-up than another. Applying the same horizons and censoring rules across models gives their comparisons a shared interpretation.

\subsection{Probability Error and Calibration}

The survival Brier score measures squared error between a predicted survival probability and the corresponding event-free status, using inverse-probability-of-censoring weighting where needed \cite{graf1999}. Integrated Brier score averages these time-specific errors over a stated evaluation window. Lower values indicate less probability prediction error. IBS combines aspects of discrimination and calibration, so it complements risk ordering while remaining distinct from a direct calibration assessment.

The window is part of the estimand. A score over early follow-up and a score over a later period evaluate different predictions, even for the same cohort and model. Training-derived horizons and common within-fold support therefore contribute to fair comparison. Within-cohort paired IBS differences are the primary predictive quantities in this thesis.

Calibration compares predicted probabilities with observed event experience. A model can preserve patient ordering while systematically overpredicting risk. Calibration intercepts, slopes and curves describe different aspects of this agreement \cite{vancalster2019,ravi60}. Austin and colleagues' work on survival calibration curves and the integrated calibration index adds a probability-scale perspective, while practical guidance emphasises the time horizon and censoring adjustment \cite{ravi60,ravi88}.

For the calibration parameterisation used here, zero intercept and unit slope are reference values. Fold means summarise average calibration behaviour, while their dispersion shows how that behaviour changes across validation samples. The combination of IBS, concordance, AUC and calibration therefore links overall prediction error to its ordering and probability components. Posterior intervals add model-conditional uncertainty to this evaluation, with empirical calibration assessed on held-out patients.

\section{Posterior Interpretation and Feature Contributions}
\label{explainable-survival-modelling}

Explanation in survival analysis depends on the quantity being explained. A regression coefficient, a hazard ratio, a risk-score contribution and a contribution to survival probability have different units and meanings. Bayesian posterior summaries address uncertainty in model parameters. Attribution methods describe how a specified fitted function uses its predictors. Combining them gives complementary views of the same model when their targets are stated clearly \cite{ravi22,ravi30,ravi61}.

\subsection{Coefficients, Hazard Ratios and Direction}

For a proportional-hazards coefficient $\beta_j$, the conditional hazard ratio is $\exp(\beta_j)$. A positive coefficient corresponds to a higher hazard per encoded predictor unit, with other predictors held fixed. Numeric standardisation makes that unit a training-set standard deviation; categorical effects refer to defined contrasts. Interpretation therefore follows the model's preprocessing rather than the original column name alone.

A posterior credible interval describes an interval containing a stated probability under the fitted posterior. Exponentiating coefficient draws produces hazard-ratio draws and preserves this uncertainty on a clinically familiar relative-rate scale \cite{ravi22}. The posterior mean hazard ratio and the exponential of the posterior mean coefficient summarise different transformations, so the selected summary should be stated.

Directional probabilities, such as $P(\beta_j>0\mid D)$, describe posterior support for an association's sign. They complement magnitude and interval width. A small, consistently positive coefficient can have high directional certainty, while a larger coefficient can have greater sign uncertainty. This separates three useful descriptions: effect magnitude, uncertainty and direction. The study reports them alongside named predictors and their encoded scales.

\subsection{From Individual Effects to Global Rankings}

A coefficient summarises a conditional effect per unit, while a contribution also depends on the predictor values to which that effect is applied. Global rankings based on posterior expected absolute contributions combine coefficient magnitude with the observed predictor distribution. Grouping encoded columns back to their source variable allows a multi-level categorical predictor to be assessed as one clinical feature. The same definition can then be used across Normal, horseshoe and continuous spike-and-slab models.

Prior-specific scale parameters provide other information. A horseshoe local scale describes its shrinkage hierarchy, and a continuous mixture weight describes its variance interpolation. These summaries enrich the description of a fitted prior, whereas a common contribution measure supports comparison across priors. Permutation importance offers another model-specific definition by measuring changes in predictive performance after perturbing a feature \cite{ravi70}.

Clinical explanation studies illustrate the use of several such tools. Alabi et al. apply LIME and SHAP in nasopharyngeal-cancer survival prediction, and survival-specific approaches adapt attribution to event-time outputs \cite{ravi28,ravi30,ravi61}. The relevant distinction is what each quantity measures. A contribution ranking explains a fitted prediction function; posterior coefficient summaries describe uncertain conditional associations under the survival model.

\subsection{SHAP Targets, Agreement and Stability}

SHAP allocates a difference between a model prediction and a reference prediction to features \cite{lundberg2017}. The selected output and background distribution define the attribution task. SurvSHAP(t) explains survival-model outputs across time, while model-agnostic survival explanations also examine time-specific predictions \cite{ravi30,ravi61}. These approaches recognise the functional nature of survival prediction, where an explanation may change with the horizon being considered.

For a linear log-risk function, exact interventional SHAP has a direct form: a feature's coefficient multiplies its deviation from the background mean \cite{shaplinear}. This provides an analytically transparent comparison for the posterior-mean Bayesian risk function. Intrinsic posterior expected absolute contributions and SHAP use the same predictor representation, but differ in posterior averaging and reference centring. Their agreement has a clear structural basis and remains an empirical quantity that can vary with the fitted effects and predictor distribution.

Stability concerns how a ranking changes across samples, seeds, computational budgets and reference backgrounds. Rank correlations, top-feature overlap and rank variability capture different aspects of that change. An average correlation can be high while a subset of comparisons is less repeatable. Evaluating the distribution of agreement and the prespecified stability checks therefore adds information about the consistency of an explanation.

The interpretation is model-based. Shapley-value explanations depend on how missing features and dependence are represented, as discussed by Kumar et al. and Ma and Tourani \cite{kumar2020,ma2020}. In the present study, the comparison uses a shared log-risk target, subjects and backgrounds, providing a defined setting for examining Bayesian intrinsic rankings, same-model SHAP and CoxPH agreement. Posterior associations and attribution rankings remain distinct from causal intervention effects.

\section{Validation and Comparative Evidence}
\label{validation-leakage-and-reproducibility.}

A predictive model includes the procedures used to prepare data and choose settings. Imputation, encoding, scaling and feature selection learn from observations, just as parameter estimation does. Nested validation separates those learning operations from the patients used for final evaluation. Inner folds select settings within outer training data, and the selected procedure is fitted before predicting the outer validation fold \cite{ravi32,ravi33}.

The same principle applies to survival-specific operations. Evaluation horizons and censoring estimates are derived from training data, and learned transformations are applied to validation patients without refitting. Applying the complete learning procedure inside each training split makes the held-out score a measure of that procedure. Shared partitions across models then support paired comparisons under the same patient mix and event conditions.

Repeated cross-validation describes how prediction changes across several training and validation samples. Training sets overlap and patients reappear in validation across repeats, creating dependence between fold results. Nadeau and Bengio examine uncertainty estimation for repeated validation and motivate corrections that reflect this overlap \cite{nadeau2003}. The mean and standard deviation of fold scores describe repeated predictive behaviour; paired contrasts and their intervals address between-model differences under the resampling design.

Benchmarks provide complementary evidence to individual clinical applications. SurvBenchmark compares methods using clinical and omics data within a common evaluation framework \cite{ravi33}. Simulation studies vary data-generating conditions more directly, while clinical applications retain the complexity of measured patients \cite{ravi37,ravi73}. Together they support evaluating performance as a relationship between the learning procedure, event information and prediction task.

Clear reporting makes these relationships assessable. Reviews of machine-learning prediction studies emphasise the need for complete descriptions of populations, predictors, validation and results, and TRIPOD+AI provides updated reporting guidance for regression and machine-learning prediction models \cite{ravi34,ravi35,collins2024}. Scientific reproducibility follows from the same detail: cohort definitions, partitions, model settings and diagnostic rules connect a reported estimate to a defined analysis. The present study uses these elements to hold the prior comparison constant across model fitting, evaluation and interpretation.

\section{Synthesis and Study Rationale}
\label{chapter-summary-and-research-gap}

\begin{table}[tbp]
\centering
\singlespacing
\small
\caption{Closest methodological studies and the focus of the present comparison.}
\label{tab:closest-work}
\begin{tabularx}{\textwidth}{@{}>{\raggedright\arraybackslash}p{2.6cm}>{\raggedright\arraybackslash}X>{\raggedright\arraybackslash}X@{}}
\toprule
Study & Focus and design & Relationship to this study\\
\midrule
SurvBenchmark \cite{ravi33} & Clinical and omics model-family benchmarking across predictive, stability and computational measures. & The present benchmark centres coefficient-prior contrasts under a common full likelihood and clinical predictor representation.\\[0.4em]
Tang et al. \cite{ravi25} & Spike-and-slab lasso Cox modelling for survival prediction and associated-gene detection. & The present comparison evaluates a continuous variance mixture alongside Normal and horseshoe priors, using shared posterior and predictive quantities.\\[0.4em]
Song et al. \cite{ravi98} & A Gaussian-and-diffused-gamma prior with rank-based Cox partial-likelihood inference, iterative conditional modes and MCMC. & The present three-prior comparison models the baseline hazard explicitly and examines matched changes in training information.\\[0.4em]
Suresh et al. \cite{ravi61} & Local LIME explanations for time-specific survival probabilities and survival curves. & The present explanation target is grouped log-risk contribution; posterior averaging, same-model SHAP agreement and sampling stability are assessed separately.\\[0.4em]
Present study & Three coefficient priors within one piecewise-exponential survival model across three clinical cohorts. & Paired prior contrasts and information perturbations connect prediction, diagnostic acceptance, posterior interpretation and explanation behaviour.\\
\bottomrule
\end{tabularx}
\end{table}

The contribution lies in this controlled empirical design. A shared likelihood, predictor matrix and validation partition make the prior contrast directly interpretable. Matched training-information changes then examine the same contrast under altered information, and the posterior/explanation analysis connects it to the fitted risk function. This yields evidence about what changes with a coefficient prior across prediction, computation and interpretation within the evaluated clinical settings.
